\subsection{FINITENESS THEOREMS}

Lemma 7. Let A be a set of unit vectors in T. If there exists a real number \(\lambda < 1\) such that \((a|a') \leqslant \lambda\) for \(a, a' \in \mathbf{A}\) and \(a \neq a'\) , then the set A is finite.

For \(a, a' \in A\) such that \(a \neq a'\) , we have

\[
\left\| a - a ^ {\prime} \right\| ^ {2} = 2 - 2 (a \mid a ^ {\prime}) \geqslant 2 - 2 \lambda .
\]

Now, the unit sphere S of T being compact, there exists a finite covering of S by sets of diameter \(<(2-2\lambda)^{1/2}\) and each of these sets contains at most one point of A, hence the lemma.

Denote by \(U(w)\) the automorphism of T associated to the affine map \(w \in W\) from E to itself. We have

\[
w (x + t) = w (x) + \operatorname{U} (w). t \text {for} t \in \mathrm{T} \text {and} x \in \mathrm{E}.
\]

This defines a homomorphism U from the group W to the orthogonal group of T; the kernel of U is the set of translations belonging to W.

THEOREM 3. (i) The set of walls of a chamber is finite.

\begin{enumerate}
\def\labelenumi{(\roman{enumi})}
\setcounter{enumi}{1}
\item
  The set of directions of hyperplanes belonging to \(\mathfrak{H}\) is finite.
\item
  The group U(W) is finite.
\end{enumerate}

Assertion (i) follows immediately from Prop. 3, (iii) and Lemma 7.

We prove (ii). Let \(\mathbf{C}\) be a chamber and \(\mathfrak{M}\) the set of its walls. The facets of \(\overline{\mathbf{C}}\) (relative to \(\mathfrak{H}\) ) are the same as those relative to \(\mathfrak{M}\) (§ 1, no. 4, Prop. 9).

Since M is finite, they are finite in number. Since a facet meets only finitely many hyperplanes belonging to H, the set of hyperplanes belonging to H and meeting \(\overline{C}\) is finite, hence so is the set A(C) of unit vectors in T orthogonal to some hyperplane belonging to H and meeting \(\overline{C}\) . Consequently, there exists a real number \(\lambda < 1\) such that \((a|a') \leqslant \lambda\) for \(a, a' \in A(C)\) and \(a \neq a'\) .

Let A be the set of unit vectors in T orthogonal to a hyperplane belonging to H. Let \(a, a' \in A\) with \(a \neq a'\) . If a and \(a'\) are parallel, then \(a = -a'\) and \((a|a') = -1\) . Otherwise, let \(H \in H\) (resp. \(H' \in H\) ) be such that a (resp. \(a'\) ) is orthogonal to H (resp. \(H'\) ). We have \(H \cap H' \neq \varnothing\) , and if \(x \in H \cap H'\) there exists an element \(w \in W\) such that \(x \in w(\overline{C})\) . The vectors \(U(w).a\) and \(U(w).a'\) then belong to A(C), we have

\[
(a \mid a ^ {\prime}) = (\mathrm{U} (w). a \mid \mathrm{U} (w). a ^ {\prime}) \leqslant \lambda ,
\]

and the set \(\mathbf{A}\) is finite by Lemma 7. Hence (ii).

Now let \(w \in W\) be such that \(\mathrm{U}(w).a = a\) for all \(a \in A\) . Then \(\mathrm{U}(w).t = t\) for all t belonging to the subspace of T generated by A. On the other hand, if \(t \in T\) is orthogonal to A, we have \(\mathrm{U}(s_{\mathrm{H}}).t = t\) for all \(H \in H\) , hence \(\mathrm{U}(w).t = t\) and finally \(\mathrm{U}(w) = 1\) . Since \(\mathrm{U}(w)(\mathrm{A}) = \mathrm{A}\) for all \(w \in W\) , we deduce that \(\mathrm{U}(\mathrm{W})\) is isomorphic to a group of permutations of the finite set A, hence (iii).

PROPOSITION 4. Let C be a chamber and N a set of walls of C. Let \(W_{N}\) be the subgroup of W generated by the orthogonal reflections with respect to the elements of N. For \(H \in N\) , denote by \(e_{H}\) the unit vector orthogonal to H on the same side of H as C. The following conditions are equivalent:

\begin{enumerate}
\def\labelenumi{(\roman{enumi})}
\item
  The group \(\mathbf{W}_{\mathfrak{N}}\) is finite.
\item
  There exists a point of \(\mathbf{E}\) invariant under every element of \(W_{\mathfrak{N}}\) .
\item
  The hyperplanes belonging to \(\mathfrak{N}\) have non-empty intersection.
\item
  The family of vectors \((e_{\mathrm{H}})_{\mathrm{H}\in \mathfrak{N}}\) is free in T.
\end{enumerate}

By property (D2) at the beginning of §3, the stabiliser in W of every point of E is finite, so (ii) implies (i).

Since the group \(W_{N}\) is generated by the set of reflections with respect to the hyperplanes belonging to N, the fixed points of \(W_{N}\) are the points of E belonging to every hyperplane \(H \in N\) , hence the equivalence of (ii) and (iii).

Assume that there exists a point \(a\) of \(E\) such that \(a \in H\) for all \(H \in \mathfrak{N}\) and let \(t \in T\) be such that \(a + t \in C\) . Since \((e_H | e_{H'}) \leqslant 0\) for \(H, H' \in \mathfrak{N}\) such that \(H \neq H'\) (Prop. 3), and since \((t | e_H) > 0\) for all \(H \in \mathfrak{N}\) , Lemma 3 (ii) implies that the \(e_H\) for \(H \in \mathfrak{N}\) are linearly independent. Consequently, (iii) implies (iv).

Suppose finally that the family \((e_{\mathrm{H}})_{\mathrm{H}\in\mathfrak{N}}\) is free. Let a be a point of E. For any hyperplane \(H\in N\) , there exists a real number \(c_{H}\) such that H consists of the points \(a+t\) of E with \((t|e_{\mathrm{H}})=c_{\mathrm{H}}\) . Since the family \((e_{\mathrm{H}})\) is free, there exists \(t\in T\) such that \((t|e_{\mathrm{H}})=c_{\mathrm{H}}\) for all \(H\in N\) , and the point \(a+t\) of E belongs to all the hyperplanes \(H\in N\) . Thus, (iv) implies (iii).

Remarks. 1) Since W is generated by reflections with respect to the walls of the chamber C, the preceding proposition gives a criterion for W to be finite. We shall return to this question in no. 9.

\begin{enumerate}
\def\labelenumi{\arabic{enumi})}
\setcounter{enumi}{1}
\tightlist
\item
  Let F be a finite-dimensional real affine space and G a group of automorphisms of F. For all \(g \in G\) , denote by \(U(g)\) the automorphism of the space of translations V of F associated to g. Assume that the image \(U(G)\) is a finite subgroup of \(\mathbf{GL}(V)\) ; then V has a scalar product invariant under \(U(G)\) (Integration, Chap. VII, §3, no. 1, Prop. 1). If, in addition, G acts properly on F when it is provided with the discrete topology, and if it is generated by reflections, we can apply to G the results of this paragraph.
\end{enumerate}

